% ============================================================
% Chapter 6: Experimental Design
%   Section: Construction of the Custom Validation Set
%     Subsection: Record Format
% Included from chapters/06-experimental-design/03-custom-validation-set/custom-validation-set.tex
% ============================================================

\subsection{Record Format}
\label{subsec:custom-format}

Every record reproduces, key for key and in the same order, the
format into which X-Fact was converted for this work
(Section~\ref{subsec:xfact}), so that the harness that scores the
system against X-Fact scores it against this set without
modification. The shared keys are listed in
Table~\ref{tab:custom-format}; the custom set appends its own keys
after them, which a reader of the X-Fact format ignores.

\begin{table}[ht]
\centering
\small
\begin{tabularx}{\linewidth}{@{}lX@{}}
\toprule
\textbf{Key} & \textbf{Content in the custom set} \\
\midrule
\multicolumn{2}{@{}l}{\emph{Shared with X-Fact}} \\
\texttt{language} & \texttt{en} or \texttt{es}. \\
\texttt{site} & Domain of the outlet in which the claim was read
(X-Fact records the fact-checking organisation here). \\
\texttt{claimant} & Who makes the claim: a person, organisation or
the outlet itself. \\
\texttt{claim} & One atomic, checkable assertion, as the article
states it. \\
\texttt{claimDate} & Publication date of the article. \\
\texttt{reviewDate} & When the label was last set. \\
\texttt{labelRaw} & The label in the annotation guide's vocabulary
(Table~\ref{tab:custom-labels}). \\
\texttt{label} & The label in the system's \texttt{Verdict}
vocabulary. \\
\texttt{referenceEvidenceLinks} & URLs of the evidence on which the
verdict rests. \\
\texttt{split} & Always \texttt{test}: no claim is ever used to
train or tune on. \\
\midrule
\multicolumn{2}{@{}l}{\emph{Added by the custom set}} \\
\texttt{id} & The record's file name (\texttt{fact001},
\texttt{fact002}, \ldots), in the order of annotation. \\
\texttt{topic} & One of the 23 classifier topics; its group is
derived. \\
\texttt{claimType} & \texttt{factual}, \texttt{numerical} or
\texttt{interpretive}. \\
\texttt{sourceTier} & Tier of the strongest evidence used
(Rule~4, Section~\ref{subsec:custom-rules}). \\
\texttt{onlyOwnSource} & Whether only the organisation the story is
about backs the claim (Rule~1). \\
\texttt{evidenceDate} & Date of the newest evidence used
(Rule~3). \\
\texttt{articleUrl} & The article the claim was taken from. \\
\texttt{annotatorNote} & The reasoning behind a close call, and the
rule that decided it. \\
\texttt{review} & The blind second label, when the record is in the
verification sample (Section~\ref{subsec:custom-review}). \\
\bottomrule
\end{tabularx}
\caption{Record format of the custom validation set. The first ten
keys are those of the X-Fact records used in this work, in the same
order.}
\label{tab:custom-format}
\end{table}

As in X-Fact, the evidence links are kept for inspecting a
disagreement, not fed to the system: the system retrieves its own
evidence from the open web at verification time, and handing it the
annotator's sources would measure the reasoning stage alone, which is
precisely the flaw for which FEVER was discarded
(Section~\ref{subsec:discard-justification}).
